\begin{frame}[label=q5a4]{5a.4 · DataFrame: nacionalidad y edad}
  \small
  \pregunta{¿Qué combinaciones de nacionalidad extranjera y rango de edad reúnen al menos 30 personas?}
  \entrada{Filtrar nacionalidades distintas de PERÚ y NO ESPECIFICA. Usar \texttt{nacionalidad}, \texttt{rango\_edad}, \texttt{cantidad}.}
  \flujo{Excluir PERÚ y\\NO ESPECIFICA}{\texttt{groupBy}\\Nacionalidad y edad}{\texttt{SUM(cantidad)}\\Total >= 30}{9 combinaciones\\Orden descendente}
  \begin{columns}[T,onlytextwidth]
    \column{0.55\textwidth}
    \textbf{Código Scala original (fragmento)}
    \codigo{sql/05a_4_nacionalidad_rangoedad}{10}{17}
    \column{0.42\textwidth}
    \textbf{Salida}\\[5pt]
    {\footnotesize
\begin{tabular}{@{}lll@{}}
\toprule
Nacionalidad & Edad & Personas\\
\midrule
Venezuela & 12--17 & 1\,009\\
Venezuela & 18--29 & 595\\
Colombia & 18--29 & 52\\
\bottomrule
\end{tabular}
    }
  \end{columns}
  \respuesta{La combinación con más personas es Venezuela y 12--17 años. Se compara el volumen registrado, sin estimar tasas por población extranjera.}
  \origen{sql}{05a_4_nacionalidad_rangoedad}
\end{frame}
